% ch21 — Camera Models and 3D Rendering
\chapter{Camera Models and 3D Rendering}
\label{ch:camera-rendering}

Before rays enter curved spacetime they must be generated from a camera model,
and after integration the resulting intensities must be assembled into a final
image. This chapter bridges classical 3D rendering with the relativistic ray
tracer.

\section{Homogeneous Coordinates and View Transformations}
\label{sec:cam-homogeneous}
\coderef{Camera.Model}{hs}

We review homogeneous coordinates in projective space and construct the
model-view-projection matrix chain that positions and orients the virtual
camera.

\section{Field of View and Perspective Projection}
\label{sec:cam-fov}

We define the field-of-view angle, derive the perspective projection matrix,
and map from clip coordinates to pixel coordinates for image generation.

\section{Lorentz Boost of the Camera Frame}
\label{sec:cam-lorentz}

We apply a Lorentz boost to the observer tetrad to model a camera in orbital
motion, accounting for aberration and the relativistic headlight effect.

\section{Supersampling and Halton Sequences}
\label{sec:cam-supersampling}

We describe anti-aliasing via supersampling with low-discrepancy Halton
sequences, achieving smooth edges and accurate sub-pixel integration.

\section{Star Catalogues\starmark{}}
\label{sec:cam-stars}
\coderef{star\_catalog}{rs}

We discuss the integration of real star-catalogue data (Hipparcos/Tycho) to
render a realistic background sky that is gravitationally lensed by the black
hole.

\section{Procedural Textures\starmark{}}
\label{sec:cam-textures}
\coderef{noise}{rs}

We implement Perlin and simplex noise to generate procedural textures for
accretion disk turbulence and background nebulae.

\paragraph{Key References.}
The camera model follows standard computer-graphics treatments
\cite{Marschner2016}; the Lorentz-boosted camera draws on \cite{James2015};
Halton sequences follow \cite{Halton1964}; procedural textures follow
\cite{Perlin1985}.
